For a law \(P\), define the finite one-sided score measure
\[
 \mu_{t,P}(E)=P_X(E\cap\mathcal A_t).
\]
If \(P\) belongs to either displayed class, \(x\in\mathcal B\), and \(r>0\), put \(s=\min\{h_0,r/2\}\).  The polynomial lower-mass property inherited from the corresponding class definition, together with \(h_0>0\) and \(c_m>0\), gives
\[
 \mu_{t,P}\{z:d(z,x)<r\}
 \ge q_{t,P}(x,s)
 \ge c_ms^{\kappa}>0. \tag{11}
\]
Thus every boundary point belongs to \(\mathcal X_{t,P}\).  The model-class Lipschitz property makes the distinguished representative \(g_{t,P}\) continuous there, so the definition of the trace gives
\[
 \theta_{t,P}(x)=g_{t,P}(x),
 \qquad t\in\{0,1\},\quad x\in\mathcal B. \tag{12}
\]
Because \(\mathcal A_0=\mathbb R^2\setminus\mathcal A_1\), sharp assignment makes \(X\in\mathcal A_t\) equivalent to \(T=t\).  Hence all side averages below have conditional observed-outcome mean \(g_{t,P}(X)\).  Consistency supplies the causal interpretation of the resulting trace contrast; the risk calculations themselves concern its identified observed-law functional.

We first prove the lower bounds by experiments contained in the stated classes.  Begin with any member of the nonempty class under consideration, retain its score law and geometry, take independent \(Z_0,Z_1\sim\mathcal N(0,1)\), and set \(Y(s)=\sigma Z_s\) for \(s\in\{0,1\}\).  Call the resulting law \(Q_0\).  Its two side regressions and traces vanish.  All design conditions are unchanged, its traces obey the envelope, and its observed residual is conditionally \(\mathcal N(0,\sigma^2)\), so \(Q_0\) remains in the same class.

For the isolated class let
\[
 h=a_0n^{-1/(\kappa+2)},
\]
where \(a_0>0\) will be fixed below.  For all sufficiently large \(n\), \(h\le h_0/3\).  The isolated thin-site property in \(\mathrm{def:isolated\text{-}thinning\text{-}class}\) and the definitions of \(\mathsf J_P\) and \(a_P\) provide \(x_h\in\mathcal B\) and \(t_h\in\{0,1\}\) such that
\[
 q_{t_h,P}(x_h,h)=a_P(x_h,h)\le C_mh^{\kappa}. \tag{13}
\]
Set
\[
 \psi_h(z)=\left(1-\frac{3d(z,x_h)}h\right)_+,
 \qquad \eta_h=\frac{Lh}{12}.
\]
If \(t_h=1\), add \(\eta_h\psi_h(X)\) to \(Y(1)\); if \(t_h=0\), subtract it from \(Y(0)\); leave the other potential outcome unchanged.  The function \(\eta_h\psi_h\) is \(L/4\)-Lipschitz, its absolute value is at most \(Lh_0/12<L\), and the score law and conditional Gaussian error law are unchanged.  Therefore the perturbed law \(Q_1\) belongs to \(\mathcal P_{\kappa}^{\mathrm{iso}}\).  By (12),
\[
 \lVert\tau_{Q_1}-\tau_{Q_0}\rVert_{\infty}\ge\eta_h. \tag{14}
\]
The two observed laws have the same score marginal.  Conditional on an assigned-side score \(z\), their Gaussian means differ by \(\eta_h\psi_h(z)\); the log density ratio for two \(\mathcal N(\mu,\sigma^2)\) laws, integrated under the shifted law, equals the squared mean difference divided by \(2\sigma^2\).  Consequently, using the support of the tent and (13),
\[
 \begin{aligned}
 \operatorname{KL}(Q_1,Q_0)
 &=\frac{\eta_h^2}{2\sigma^2}
   \int\psi_h(z)^2\,d\mu_{t_h,P}(z)\\
 &\le\frac{C_mL^2}{288\sigma^2}h^{\kappa+2}. 
 \end{aligned} \tag{15}
\]
Factorization of the likelihood under independent sampling makes relative entropy additive, whence
\[
 \operatorname{KL}(Q_1^{\otimes n},Q_0^{\otimes n})
 \le C a_0^{\kappa+2}.
\]
Choose \(a_0\) so that the last bound is at most \(1/8\).  The two-law clause of \(\mathrm{lem:finite\text{-}information\text{-}testing}\), with the supremum metric and \(2\delta=\eta_h\), now gives
\[
 R_n(\mathcal P_{\kappa}^{\mathrm{iso}})
 \ge c n^{-1/(\kappa+2)}
 =c r_{n,\kappa}^{\mathrm{iso}}. \tag{16}
\]

For the pervasive class let
\[
 h=a_0\left(\frac{\log n}{n}\right)^{1/(\kappa+2)}.
\]
For large \(n\), \(h\le h_0/3\).  The pervasive thin-site property in \(\mathrm{def:pervasive\text{-}thinning\text{-}class}\) supplies an \(h\)-separated set \(x_1,\ldots,x_M\) such that
\[
 M\ge c_Jh^{-1},
 \qquad a_P(x_j,h)\le C_mh^{\kappa}\quad(1\le j\le M). \tag{17}
\]
Its finiteness follows from global boundary entropy.  At each \(x_j\), choose a side attaining the minimum in \(a_P(x_j,h)\), and perturb that side by a radius-\(h/3\), height-\(Lh/12\) tent, upward on side one and downward on side zero.  The membership argument used above gives laws \(Q_1,\ldots,Q_M\in\mathcal P_{\kappa}^{\mathrm{perv}}\).  At \(x_j\), the target under \(Q_j\) differs from the zero target by \(Lh/12\), whereas every other tent vanishes there because the centers are \(h\)-separated.  Thus the target curves are pairwise separated by at least \(Lh/12\) in supremum norm.  The calculation in (15) gives
\[
 \max_{1\le j\le M}
 \operatorname{KL}(Q_j^{\otimes n},Q_0^{\otimes n})
 \le Cnh^{\kappa+2}
 =Ca_0^{\kappa+2}\log n. \tag{18}
\]
For fixed \(\kappa\), (17) and the chosen value of \(h\) imply \(\log M\ge c_{\kappa}\log n\) eventually.  Choose \(a_0\) still smaller if necessary so that (18) is at most \(a\log M\), where \(a<1/8\) is the constant in \(\mathrm{lem:finite\text{-}information\text{-}testing}\).  Its multiple-law clause yields
\[
 R_n(\mathcal P_{\kappa}^{\mathrm{perv}})
 \ge c\left(\frac{\log n}{n}\right)^{1/(\kappa+2)}
 =c r_{n,\kappa}^{\mathrm{perv}}. \tag{19}
\]
This preserves the exact-Gaussian witnesses and Gaussian Kullback--Leibler calculations; \(\mathrm{lem:fano\text{-}tent\text{-}validity}\) separately certifies explicit circular families with the same properties.

We next construct data-driven upper-bound witnesses for the infimum defining \(R_n\).  Let \(s_n\) denote either \(r_{n,\kappa}^{\mathrm{iso}}\) or \(r_{n,\kappa}^{\mathrm{perv}}\).  By \(\mathrm{def:dyadic\text{-}grid}\), for all sufficiently large \(n\) there is \(b=b_n\in\mathcal H_n\) such that
\[
 s_n\le b<2s_n,
 \qquad b/2\in\mathcal H_n,
 \qquad b\le h_0/2. \tag{20}
\]
Choose a deterministic maximal \(b/2\)-separated set \(z_1,\ldots,z_M\) in \(\mathcal B\).  Maximality makes it a \(b/2\)-net, compactness makes it finite, and global boundary entropy gives \(M\le2C_{\mathcal B}b^{-1}\).  Let \(j(x)\) be the least index satisfying \(d(x,z_{j(x)})\le b/2\), and define
\[
 K_{t,j}=\sum_{i=1}^n
 \mathbf1\{X_i\in\mathcal A_t, d(X_i,z_j)\le b\},
\]
\[
 \overline g_{t,j}
 =\frac{\sum_{i=1}^nY_i
 \mathbf1\{X_i\in\mathcal A_t, d(X_i,z_j)\le b\}}
 {K_{t,j}\vee1},
\]
and the clipped estimator
\[
 \widetilde\tau_{n,b}(x)
 =\Pi_{[-2L,2L]}\{\overline g_{1,j(x)}-\overline g_{0,j(x)}\}.
 \tag{21}
\]
This is a single estimator measurable in the sample and known assignment geometry.  Since the bounded-trace property gives \(\tau_P(x)\in[-2L,2L]\), projection cannot increase its pointwise error.  On \(K_{t,j(x)}>0\), every included score is at distance at most \(3b/2\) from \(x\).  Equations (11)--(12) and the class Lipschitz property therefore give
\[
 \left|\mathbb E_P(\overline g_{t,j(x)}\mid X_1,\ldots,X_n)
       -\theta_{t,P}(x)\right|
 \le\frac{3Lb}{2}. \tag{22}
\]

We will also use an elementary count bound.  If \(N\sim\operatorname{Bin}(n,q)\), Markov's inequality applied to \(2^{-N}\), followed by \(1-y\le e^{-y}\), gives
\[
 \Pr(N\le nq/2)
 \le2^{nq/2}(1-q/2)^n
 \le e^{-nq/8}. \tag{23}
\]

Consider first the pervasive rate and take \(s_n=r_{n,\kappa}^{\mathrm{perv}}\) in (20).  The polynomial lower-mass property and (23), followed by a union bound over the two sides and the net, give an event \(H_b\) on which every \(K_{t,j}\ge nc_mb^{\kappa}/2\) and
\[
 \sup_{P\in\mathcal P_{\kappa}^{\mathrm{perv}}}P(H_b^c)
 \le Cb^{-1}e^{-cn b^{\kappa}}. \tag{24}
\]
For fixed scores, every nonempty index set defining \(\overline g_{t,j}\) is one of the disk patterns in \(\mathrm{lem:disk\text{-}pattern\text{-}sub\text{-}gaussian\text{-}control}\), because \(z_j\in\mathcal B\) and \(b\in\mathcal H_n\).  That lemma, together with
\[
 |\mathcal H_n|\le C\log(en),
 \qquad
 \sum_{r=0}^3\binom nr\le C(1+n^3),
\]
implies that the conditional expectation of the maximum normalized residual sum over all the averages in (21) is at most \(C\sqrt{\log(en)}\).  Consequently, on \(H_b\),
\[
 \mathbb E_P\left[
 \max_{t,j}\left|
 \overline g_{t,j}
 -\mathbb E_P(\overline g_{t,j}\mid X_1,\ldots,X_n)
 \right|
 \ \middle|\ X_1,\ldots,X_n\right]
 \le C\sigma\sqrt{\frac{\log(en)}{n c_mb^{\kappa}}}. \tag{25}
\]
Because (20) implies \(nb^{\kappa+2}\asymp\log n\), the right side is at most \(C_{\kappa}b\).  On \(H_b^c\), clipping bounds the loss by \(4L\).  Combining (22), (24), and (25), and observing that the exponential remainder in (24) is \(o(b)\), yields
\[
 \sup_{P\in\mathcal P_{\kappa}^{\mathrm{perv}}}
 \mathbb E_P\lVert\widetilde\tau_{n,b}-\tau_P\rVert_{\infty}
 \le C_{\kappa}b
 \le C_{\kappa}r_{n,\kappa}^{\mathrm{perv}}. \tag{26}
\]

For the isolated rate take \(s_n=r_{n,\kappa}^{\mathrm{iso}}\) in (20), so \(nb^{\kappa}\asymp b^{-2}\).  Fix a side, choose \(D\ge C_m\), set
\[
 u_\ell=\min\{1,D2^\ell b^{\kappa}\},
\]
and put a net center in group zero when \(q_{t,P}(z_j,b)\le u_0\), and in group \(\ell\ge1\) when
\[
 u_{\ell-1}<q_{t,P}(z_j,b)\le u_\ell.
\]
Only the first index with \(u_\ell=1\) can be a terminal group.  Every center in group \(\ell\) satisfies
\[
 a_P(z_j,b/2)
 \le q_{t,P}(z_j,b/2)
 \le q_{t,P}(z_j,b)
 \le u_\ell.
\]
The centers are \(b/2\)-separated.  Applying the isolated upper profile from \(\mathrm{def:isolated\text{-}thinning\text{-}class}\) at radius \(b/2\), including at the terminal group where \(u_\ell=1\), gives
\[
 M_\ell\le C_{\kappa}\{1+2^{\ell/(\kappa-2)}\}. \tag{27}
\]
In group zero, polynomial lower mass makes the mean count at least \(c_mnb^{\kappa}\ge c b^{-2}\); in each nonempty group \(\ell\ge1\), it is at least \(D2^{\ell-1}nb^{\kappa}\ge c2^\ell b^{-2}\).  Equations (23) and (27), summed over the groups and both sides, give an event \(H_b\) on which every group-\(\ell\) count is at least \(c2^\ell b^{-2}\) and
\[
 \begin{aligned}
 \sup_{P\in\mathcal P_{\kappa}^{\mathrm{iso}}}P(H_b^c)
 &\le C_{\kappa}\sum_{\ell\ge0}
 \{1+2^{\ell/(\kappa-2)}\}e^{-c2^\ell/b^2}\\
 &\le C_{\kappa}e^{-c_{\kappa}/b^2}. 
 \end{aligned} \tag{28}
\]
The last inequality follows because, for sufficiently small \(b\), the ratio of consecutive summands is bounded above by a fixed number smaller than one.

Conditionally on the scores, independent sampling and the exact conditional Gaussian error property inherited from the isolated class make each nonempty residual average Gaussian with variance \(\sigma^2/K_{t,j}\).  Hence, on \(H_b\), a two-sided Gaussian tail bound and a union bound within each group imply, for \(z>0\),
\[
 \Pr_P\left\{
 \max_{t,j}\left|
 \overline g_{t,j}
 -\mathbb E_P(\overline g_{t,j}\mid X_1,\ldots,X_n)
 \right|>zb
 \ \middle|\ X_1,\ldots,X_n\right\}
 \le C_{\kappa}\sum_{\ell\ge0}
 \{1+2^{\ell/(\kappa-2)}\}e^{-c z^22^\ell}. \tag{29}
\]
For all sufficiently large fixed \(z\), the series in (29) is at most \(C_{\kappa}e^{-c_{\kappa}z^2}\).  Integrating this tail, and using the trivial probability bound one below that fixed threshold, gives a conditional expected stochastic maximum at most \(C_{\kappa}b\).  Equations (22), (28), and the \(4L\) clipped-loss bound on \(H_b^c\) now yield
\[
 \sup_{P\in\mathcal P_{\kappa}^{\mathrm{iso}}}
 \mathbb E_P\lVert\widetilde\tau_{n,b}-\tau_P\rVert_{\infty}
 \le C_{\kappa}b
 \le C_{\kappa}r_{n,\kappa}^{\mathrm{iso}}. \tag{30}
\]

The estimator (21) is an admissible data-driven witness for the common infimum in \(\mathrm{def:expected\text{-}sup\text{-}risk}\).  Thus (26) and (30) prove the two upper bounds, while (16) and (19) prove the matching lower bounds.  The explicit nonempty circular specifications are supplied by \(\mathrm{lem:circular\text{-}witness\text{-}membership}\).  Finally, the two classes have the same lower modulus \(c_mh^{\kappa}\), whereas
\[
 \frac{r_{n,\kappa}^{\mathrm{perv}}}
 {r_{n,\kappa}^{\mathrm{iso}}}
 =(\log n)^{1/(\kappa+2)}\longrightarrow\infty.
\]
Therefore the scalar modulus alone cannot determine expected-sup minimax difficulty, and the asserted thickness--entropy separation follows.